\begin{center}
	\large{\textbf{A THREE-QUBIT ENTANGLEMENT CLASSIFICATION STUDY USING CONCURRENCE FILL}}
\end{center}
%\begin{flushleft}
\begin{minipage}{4cm}
	\begin{align*}
		&\text{Name} \qquad &&: \text{Tamara Pingki} \\
		&\text{NRP} \qquad &&: \text{6001241001} \\
		&\text{Department} \qquad &&: \text{Fisika} \\
		&\text{Supervisor} \qquad &&: \text{Prof. Agus Purwanto, D.Sc.} \\
	\end{align*}
\end{minipage} \\
\begin{center}
	\large\textbf{ABSTRACT} \\  
\end{center}
True multipartite entanglement is a fundamental phenomenon in quantum mechanics that plays a crucial role in quantum information theory and applications. However, the characterization and measurement of true multipartite entanglement still face theoretical challenges, particularly in multipartite systems and mixed states. One emerging approach is the geometric approach, which utilizes the structure of quantum state space to quantify entanglement. This study examines concurrency fill as a geometric measure of true multipartite entanglement in three-qubit systems. The analysis focuses on the mathematical formulation of concurrency fill and its relationship to three tangle and partial tangle, thus enabling the classification of entanglement into GHZ and W-class states. In addition, the application of concurrency fill to certain mixed states, specifically rank-2 mixtures of GHZ and W-class states, is also discussed to examine the emerging entanglement structure. This study is conducted theoretically and analytically, without involving numerical simulations or experimental verification.

\noindent
\begin{minipage}{1cm}
	\begin{align*}
		\text{\textbf{Kata kunci}}:\hspace{5mm}&\text{Entanglement, Three-Qubit, Concurrence Fill}% \\ &\text{} \\&\text{}
	\end{align*}
\end{minipage}


\newpage
\thispagestyle{empty}
\begin{center}
	\topskip0pt
	\vspace*{\fill}
	\textit{\textbf{"Halaman ini sengaja dikosongkan"}}
	\vspace*{\fill}
\end{center}
\pagebreak